\chapter{Projekt systemu i metodyka badawcza}
\label{ch:projekt}

\section{Założenia i wymagania systemowe}
\todo{Wymagania funkcjonalne (ocena transakcji, zarządzanie modelami, metryki, historia) i niefunkcjonalne (opóźnienie, modułowość); tabela.}

\section{Architektura klient--serwer i dobór narzędzi}
\todo{Diagram: frontend (React), API (FastAPI), moduł ML (scikit-learn, XGBoost), baza (SQL).}
Architektura jest modułowa: źródła transakcji (symulator, odtwarzanie zbioru, import CSV, szkielet konektora bankowego) dostarczają ujednoliconych obiektów \texttt{Transaction}, a moduł oceny mapuje je na cechy właściwe dla schematu danego modelu. Dzięki temu dodanie nowego źródła lub modelu nie wymaga zmian pozostałych części. \todo{rozwinąć po ustabilizowaniu kodu}

\section{Model relacyjnej bazy danych}
\todo{Schemat tabel (historia ocen, modele) i uzasadnienie SQLite/PostgreSQL.}

\section{Charakterystyka i podział zbioru danych}
Zbiór \emph{Credit Card Fraud Detection} \cite{kaggle_ulb} zawiera 284\,807 transakcji kartowych z dwóch dni, w tym 492 oszustwa (0{,}172\%). Cechy V1--V28 są wynikiem transformacji PCA (ze względu na poufność), a ponadto dostępne są czas od pierwszej transakcji, kwota i etykieta klasy. \todo{cytowanie zbioru zgodnie z opisem na stronie, liczność klas w podzbiorach, uwagi o ograniczeniach.}

Dane dzielono chronologicznie na zbiory treningowy, walidacyjny i testowy (70/15/15), tak aby testowanie odpowiadało predykcji „przyszłości” względem uczenia. \todo{rozkład oszustw w podzbiorach, opis wariantu odtwarzającego ustawienie z \cite{niu2019}}
